\documentclass[ruledheader,noindentfirst,anapcustomindent,abntfigtabnum,tocpage=plain]{abnt}
\usepackage{amsmath, amssymb, amsthm, verbatim, amsfonts, amstext}
%\usepackage[latin1]{inputenc}
\usepackage[brazilian]{babel}
\usepackage[utf8]{inputenc}
\usepackage[T1]{fontenc}
\usepackage{dropping}
\usepackage{graphicx}
\usepackage[hang,small,bf]{caption}
\usepackage[abnt-etal-list=0,abnt-etal-text=it,abnt-and-type=&,abnt-emphasize=bf,abnt-full-initials=yes,alf,bibjustif]{abntcite}
\usepackage{fancyhdr}
\usepackage{makeidx}
\usepackage[none]{hyphenat}
\usepackage{color}
\usepackage{subfig}
\usepackage{algorithms}
\usepackage{algorithmic}
\usepackage{mdwlist}
\usepackage{bm}
\usepackage[titletoc,title]{appendix}
\usepackage{ltxtable}
\usepackage{longtable}
\usepackage{supertabular}
\usepackage{indentfirst}
\usepackage{color}
\usepackage{icomma}

\sloppy


%
%Tradução do pacote Algorithm para portugues
%
\renewcommand{\algorithmicrequire}{\textbf{Entrada:}}
\renewcommand{\algorithmicensure}{\textbf{Saída:}}
\renewcommand{\algorithmicend}{\textbf{fim}}
\renewcommand{\algorithmicif}{\textbf{se}}
\renewcommand{\algorithmicthen}{\textbf{então}}
\renewcommand{\algorithmicelse}{\textbf{senão}}
\renewcommand{\algorithmicelsif}{\algorithmicelse \, \algorithmicif}
\renewcommand{\algorithmicendif}{\algorithmicend \, \algorithmicif}
\renewcommand{\algorithmicfor}{\textbf{para}}
\renewcommand{\algorithmicforall}{\textbf{para todo}}
\renewcommand{\algorithmicdo}{\textbf{fazer}}
\renewcommand{\algorithmicendfor}{\algorithmicend \, \algorithmicfor}
\renewcommand{\algorithmicwhile}{\textbf{enquanto}}
\renewcommand{\algorithmicendwhile}{\algorithmicend \, \algorithmicwhile}
\renewcommand{\algorithmicloop}{\textbf{laço}}
\renewcommand{\algorithmicendloop}{\algorithmicend \, \algorithmicloop}
\renewcommand{\algorithmicrepeat}{\textbf{repetir}}
\renewcommand{\algorithmicuntil}{\textbf{até}}
\renewcommand{\algorithmiccomment}[1]{\{#1\}}
\renewcommand{\listalgorithmname}{Lista de Algoritmos}
\floatname{algorithm}{Algoritmo}
%%%%%%%%%%%%%%%%%%%%%%%%%%%%%%%%%%%%%%%%%%%%%%%%%%%%%%%%%%%%%%%%%%%%%%%%%%%%%%%%%%%

\makeindex

%%%% O arquivo modelosCAP.tex possui as definições para ciação do estilo de capítulo (fonte de título, barras horizontais, etc.)
% ele não gera texto de saída, é um arquivo de configuração somente
%
%Estilo de formatação de capítulos

\makeatletter
\newcommand{\thechapterwords}
{ \ifcase \thechapter\or 1\or 2\or 3\or 4\or 5\or6\or 7\or 8\or 9\or 10\or 11\fi}

\def\@makechapterhead#1{%
\vspace*{10\p@}%
{\parindent \z@  \reset@font

\scshape \@chapapp{} \thechapterwords
\quad %
\par\nobreak
\vspace*{10\p@}%
\interlinepenalty\@M
\hrule
\vspace*{10\p@}%
\Huge \bfseries #1\par\nobreak
\par
\vspace*{10\p@}%
\hrule
\vskip 40\p@
}}
\def\@makeschapterhead#1{%
\vspace*{10\p@}%
{\parindent \z@ \centering \reset@font
\par\nobreak
\vspace*{10\p@}%
\interlinepenalty\@M
\hrule
\vspace*{10\p@}%
\Huge \bfseries #1\par\nobreak
\par
\vspace*{10\p@}%
\hrule
\vskip 40\p@
%\vskip 100\p@
}}
%%%%%%%%%%%%%%%%%%%%%%%%%%%%%%%%%%%%%%%%%%%%%%%FIM DO PREAMBULO%%%%%%%%%%%%%%%%%%%%%%%%%%%%%%%%%%%%%%%%%%%%%%%%%%%%%%%%%%%%%%%%%%


\begin{document}

%%%%% IMPORTANTE: ALTERA O TEXTO ENTRE ARIAL E TIMES NEW ROMAN (ALTERNAR OS COMENTÁRIOS)
%
%%%%%%%%%%%%%%%%%%%%%PARA UTILIZAR ARIAL%%%%%%%%%%%%%%%%%%%%%%%
%
\fontfamily{phv}                    %fonte Arial
\renewcommand{\rmdefault}{phv}      %
%
%%%%%%%%%%%%%%%%%%%%%PARA UTILIZAR TIMES%%%%%%%%%%%%%%%%%%%%%%%
%
%\fontfamily{ptm}               %fonte Times
%\renewcommand{\rmdefault}{ptm} %
%
%%%%%%%%%%%%%%%%%%%%%%%%%%%%%%%%%%%%%%%%%%%%%%%%%%%%%%%%%%%%%%%

%%%%%%%%%%%%%Arquivos .tex com os elementos pré-textuais
%
\thispagestyle{empty}

\vfill
 \begin{center}
    \begin{figure}[t]
     \centering
            \includegraphics[width=5cm]{figures/IF_logo.eps}\\[-0.1in]
     \end{figure}

    {\large\bfseries INSTITUTO FEDERAL DE EDUCAÇÃO, CIÊNCIA E TECNOLOGIA DO PARÁ} \\
    {\large\bfseries PRÓ-REITORIA DE ENSINO} \\
    {\large\bfseries COORDENADORIA DE TELEMÁTICA DO CAMPUS MARACANAÚ}  \\ 
    {\large\bfseries BACHARELADO EM AGRONOMIA}  \\ 

    \vspace*{1in}
    \begin{large} \bfseries LUCAS PALHETA SAMPAIO \end{large}\\[0.4in]

    \vspace*{4cm}
    \noindent \\
    \large\bfseries{FERRAMENTA EM AMBIENTE WEB PARA O DIMENSIONAMENTO DE SISTEMAS DE IRRIGAÇÃO} \\
    \vfill
    \large\bfseries{ CASTANHAL \\ 2016}
\end{center}

\normalsize
\begin{titlepage}
\vfill
\begin{center}

    {\large LUCAS PALHETA SAMPAIO\\}
    \vspace{2cm}
    {\Large \textsc{FERRAMENTA EM AMBIENTE WEB PARA DIMENSIONAMENTO DE SISTEMAS DE IRRIGAÇÃO}\\}
    \vspace{1cm}
    \hspace{.45\linewidth}
    \begin{minipage}{.50\linewidth}

            Trabalho de conclusão de curso apresentado ao Instituto Federal de Educação, Ciência e Tecnologia do Pará - Campus Castanhal, como requisito parcial para a obtenção de Grau em Bacharelado em Agronomia.
            

            \vspace{0.5 cm}

            Área de pesquisa: Irrigação

            \vspace{0.5 cm}

            Orientador: Luiz Nery Rodrigues
    
    \end{minipage}

    \vspace{2cm}
    \vfill
    {\large Castanhal \\ 2016}
\end{center}

\end{titlepage}
\begin{folhadeaprovacao}
\setlength{\ABNTsignthickness}{0.2pt}
\setlength{\ABNTsignskip}{1.7cm}

\begin{center}
\includegraphics[width=2.5cm]{figures/brasao_republica.eps}\\
%\includegraphics[width=5cm]{figures/IF_logo.eps}\\ %outros brasões
%\includegraphics[width=5cm]{figures/IF_logo2.eps}\\%outros brasões

            {INSTITUTO FEDERAL DE EDUCAÇÃO, CIÊNCIA E TECNOLOGIA DO CEARÁ} \\
            {COORDENAÇÃO DE PÓS-GRADUAÇÃO EM ENGENHARIA DE TELECOMUNICAÇÕES}  \\

    \vspace{1.5cm}
                                    {FELIPE MARCEL DE QUEIROZ SANTOS}\\
    \bfseries{}
\end{center}

Esta Monografia foi julgada adequada para a obten\c{c}\~{a}o do Grau de Bacharel em Ciência da Computação, sendo aprovada pela Coordenadoria de Telemática e pela Coordenadoria do curso de Bacharelado em Ciência da Computação do Campus Maracanaú do Instituto Federal de Educação, Ciência e Tecnologia do Ceará e pela banca examinadora:

    \vspace{0.15cm}
    \assinatura{Orientador: Prof. Dr. Amauri \\ Instituto Federal do Ceará - IFCE}
    \assinatura{Prof. Dr. Huguinho \\ Instituto Federal do Ceará - IFCE}
    \assinatura{Prof. Dr. Zezinho \\ Instituto Federal do Ceará - IFCE}
    \assinatura{Prof. Dr. Luizinho \\ Instituto Federal do Ceará - IFCE}
    \vspace{0.15cm}%\vfill

    \begin{center}
        Fortaleza, 06 de Abril de 2013
    \end{center}
\end{folhadeaprovacao}
\vspace*{15cm}

\hfill Dedico este trabalho ...\\
\chapter*{Agradecimentos}
%\thispagestyle{empty}


\begin{flushright}
\begin{minipage}[r]{10cm}
\vspace{18cm}
``A mente que se abre a uma nova idéia jamais voltará ao seu tamanho original''.
\begin{flushright}
Albert Einstein
\end{flushright}
\end{minipage}
\end{flushright}
\pagestyle{plain}%%%%% Utilizar ESTILO PLAIN AQUI%%%%%%%
\chapter*{Resumo}

\noindent Este trabalho apresenta...
\chapter*{Abstract}


\noindent This work presents...

%%%Comandos para criação automática das listas
%
\tableofcontents
\listoffigures
\listoftables

%%%Comandos para criar outras listas não suportadas pelo pacote ABNTex%%%
%
\pretextualchapter{Lista de Símbolos}
\begin{basedescript}{\desclabelstyle{\pushlabel}\desclabelwidth{6em}}
\item[$Z$] variavel aleatoria%
\item[$\mathbb{R}$] conjunto dos números reais%
\item[$t$] tempo contínuo%
\item[$n$] tempo discreto%
\item[$f(z)$] função densidade de probabilidade%
\item[$F(z)$] função de distribuição acumulada%
\item[$\sigma$] desvio padrão%
\item[$\mu$] média ou esperança matemática%
\item[$|\cdot|$] operador magnitude%
\item[$\nabla$] operador gradiente%
\end{basedescript}
\newpage

\pretextualchapter{Lista de Abreviacoes}
\begin{basedescript}{\desclabelstyle{\pushlabel}\desclabelwidth{6em}}
\item[{fdp}] Função densidade de probabilidade%
\item[{fda}] Função de distribuição acumulada%
\item[{EMQ}] Erro médio quadrático%
\end{basedescript}
\newpage
%%%%%%%%%%%%%%%%%%%%%%%%%%%%%%%%%%%%%%%%%%%%%%%%%%%%%%%%%%%%%%%%%%%%

%Capítulos passam a ter páginas numeradas
%
\pagestyle{fancy}

%resseta os contadores de capítulo e seção
%
\renewcommand{\chaptermark}[1]{\markboth{#1}{}}
\renewcommand{\sectionmark}[1]{\markright{\thesection\ #1}}

%%%%%%%%%%%%%%NÃO LEMBRO O QUE FAZ, APARENTEMENTE NADA, TESTAR DEPOIS
%\fancyhf{}%
%\fancyhead[RO,LE]{\large\slshape\thepage}%
%\fancyhead[CE]{\large\slshape\leftmark}%
%\fancyhead[CO]{\large\slshape\rightmark}%


%%% Outros arquivos .tex. É acoselhável utilizar vários arquivos, pelo menos um por capítulo
\chapter{Introdução}\label{CAP:introducao}
%\thispagestyle{empty}

Este documento consiste de um modelo basico para a producao de documentos academicos, seguindo as normas ABNT. 

Nao e abordado o estudo do LaTex neste template. Sugerimos a leitura do texto em \citeonline{Oetiker:1995}. O uso do LaTex e aconselhavel devido a sua qualidade grafica, facil referenciacao, criacao de listas, indices, referencias bibliograficas e escrita matematica profissional. Porem, nao e obrigatorio o uso deste template, apenas as orientacoes de formatação segundo as normas ABNT devem ser obrigatoriamente seguidas.

Uma observação em particular é a de que, no pacote ABNTex, as referências diretas devem utilizar o comando ``citeonline''. Referências indiretas utilizam o comando ``cite''.

Exemplo de citacao direta: Uma otima fonte de estudo para compreender o LaTex e apresentada por \citeonline{Oetiker:1995}. 

Exemplo de citação indireta: Existem boas fontes de pesquisa para entendimento do LaTex \cite{Oetiker:1995}, estas incluem documentação online disponível na web.

\section{Motivação e objetivos}


 
\section{Contribuicoes}




\section{Producao cientifica}


\section{Organizacao da tese}

\noindent \textbf{Capitulo \ref{CAP2}}: descricao...

\noindent \textbf{Capitulo \ref{CAP3}}: descricaoo...

\noindent \textbf{Capitulo \ref{CAP4}}: descricao...

\noindent \textbf{Capitulo \ref{CAP5}}: descricao...
\chapter{Panorama Geral de Irrigação}
\label{CAP2}



\section{Conjuntura Mundial}\label{Sub:equa}

A irrigação, utilizada como técnica primordial para os cultivos em áreas com deficit hídrico elevado tem anualmente expandido sua área global, dados do FAOSTAT (2016) mostram que a área irrigada no mundo no ano de 2010 estava na casa dos 320 milhões de hectares, tendo um aumento de 5 milhões de hectares até o ano de 2013.
A mesma instituição afirma que nos últimos 10 anos o país teve um crescimento de 800 mil hectares irrigados, com acréscimo média de 200 mil por ano de 2006 até 2010, estabilizando em 5400 mil hectares até 2013.  

A Agencia Nacional de Água (ANA), órgão que monitora os recursos hídricos do país, informou em 2015 que a demanda conjuntiva de água chegou a 2.275 m cubico/s, tendo como maior contribuinte desse índice o setor de irrigação, detentor da parcela de 55 (porcento) sendo que o segundo maior consumidor de água é o abastecimento humano urbano com apenas 22 (porcento), contudo .

A área brasileira irrigada no ano de 2014 foi estimada em 6,11 milhões de hectares  ou 21 (porcento) do potencial nacional, que corresponde a 29,6 milhões de hectares, contudo, observa-se expressivo aumento da agricultura irrigada no Brasil, crescendo sempre a taxas superiores às do crescimento da área plantada total.

Investimentos em irrigação resultam em aumento substancial da produtividade e do valor da produção agrícola, diminuindo a necessidade de expansão em áreas ocupadas por outros usos e coberturas (pastagens ou matas nativas, por exemplo). Aplicando boas práticas de manejo do solo e da água, irrigantes alcançam efciências de uso dos recursos hídricos superiores a 90 (porcento). \cite{ana2015}.

\section{Uso Indevido da Água na Irrigação}
O crescimento populacional tem gerado demanda a cada ano por água, em consequência, a disponibilidade desse recurso vem tornando-se menor gradativamente, e o reflexo vem sendo observado não só no na escassez do elemento em si como na expansão das fronteiras agrícolas e à degradação do meio ambiente. Sendo a água um recurso indispensável à vida, é de fundamental importância a discussão das relações entre o homem e a água, uma vez que a sobrevivência das gerações futuras depende diretamente das decisões que hoje estão sendo tomadas \cite{lima1999uso}.

A escolha do sistema mais adequado depende de diversos fatores, \cite{de2006avaliaccao} enumera a topografia e o tipo de solo, a fonte de água (localização, vazão, qualidade), o sistema de plantio e o custo do equipamento e de operação. Contudo, no mesmo trabalho afirma que os sistemas que inicialmente tiveram um inadequado dimensionamento hidráulico, tiveram sua uniformidade de aplicação deficitária, causando decréscimo na produtividade devido ao uso irracional da água.


\section{Efeitos da Irrigação na Produção}

Em comparação com áreas não irrigadas a produção por hectare de culturas sob regime de irrigação, demonstra acrécimos em diversas áreas de cultivo. \cite{sanches2013girassol} desenvolvendo trabalhos com girassol, obteve altas significativas em áreas irrigadas, alcançado taxas de 62 (porcento) a mais que em áreas sem regime de irrigação. 

No cenário da pecuária, os estudos de \cite{sanches2013tifton} com capim tifton 85   sobresemeado com aveia, demonstraram índices mais elevados de matéria seca em kg/ha a partir do segundo siclo de pastejo. Os valores alcançaram a faixa de 82 (porcento) a mais de matéria seca nas parcelas irrigadas. O nível de proteína bruta também verificado apresentou acréscimos significativos.

Mesmo sendo a técnica de produção agrícola com utilização de um volume de água demasiadamente grande, essa é uma ação necessária, pois a aplicação de água nas culturas aumenta a eficiência de uso de outros insumos, como fertilizantes, por exemplo, garante a produção na entressafra em regiões áridas ou de regime pluviométrico inconstante, além de oferecer segurança durante os veranicos \cite{de2008desenvolvimento}.

\section{Relação Água-Solo-planta}

Para se obter uma utilização adequada de dos recursos hídricos na atividade de irrigação, é imprescindível que o planejamento venha a abranger as inter-relações entre os fatores solo, água, planta e a atmosfera.

A quantidade de água necessária para a irrigação é um dos principais parâmetros a ser determinada em um correto planejamento de dimensionamento do sistema. Se superestimada gera consequências sistemas de irrigação superdimensionados, encarecendo custo de irrigação por unidade de área, aplicação em excesso, com consequente reflexo nas características físico-químicas do solo, como lixiviação dos nutrientes e salinização.

Em sistemas sub-estimados, por não suprir as necessidades hídricas ideais, a cultura apresentará desenvolvimento a quem do esperado, negligenciando assim o potencial de produção, e ganhos econômicos.

A água necessária é definida por \cite{bernardo2006} como sendo a quantidade de água requerida pela cultura, em determinado período de tempo, de modo a não limitar seu crescimento e sua produção, nas condições climáticas locais, ou seja, a quantidade de água necessária para atender à evapotranspiração e à lixiviação dos sais do solo. Assim fica evidente que dois parâmetros influenciam de forma ativa na quantidade de água requerida, a evapotranspiração e a lixiviação.

Por definição a evapotranspiração consiste na somatória da evaporação do solo, e da transpiração ocorrida na planta durante determinado período. Para que tal processo ocorra é necessário que haja temperatura e umidade do ar, gerando a energia necessária para evaporação da água. Observa-se que o a atmosfera atua como dreno do solo, retirando a água aderida nas partículas.

\section{Tipos de Sistemas}

\subsection{Sistema de Aspersão}

O sistema de aspersão convencional é o sistema básico de aspersão, esse sistema é caracterizado pela utilização de uma rede de tubulações móveis de engate rápido, ou fixo e enterrado, irrigando áreas pequenas ou médias.

Esse modelo de sistema pode ser classificado segundo sua mobilidade, apresentando possibilidade de sistemas portáteis onde toda a tubulação é móvel, inclusive a motobomba; pode ser observado também como sistemas semiportáteis, no qual a linha principal é fixa e as laterais são móveis; no sistema fixo até as linhas laterais são fixas; e no sistema em malha as tubulações são fixas porém enterradas; uma configuração de sistemas de aspersão com grandes vazões é o sistema de canhão hidráulico, muito utilizado na irrigação de pastagens.

Os sistemas de irrigação por aspersão apresentam aspersores com maior vazão, sendo necessário maior cuidado no momento do dimensionamento hidráulico para que essa vazão não ultrapasse a velocidade de infiltração básica, ocasionando (esqueci o termo técnico para lavagem superficial do solo).

Em geral, os sistemas de aspersão são constituídos por um conjunto moto bomba, tubulações de sucção, recalque, principal, com linhas laterais, há também a possibilidade de apresentar, acoplada a linha principal, linhas de derivação as quais tem a função de distribuição da água para as linhas laterais. Apresenta também acessórios como válvulas e reguladores de pressão, além de aspersores.

\subsection{Sistema de Gotejamento}

\chapter{Métodologia}\label{CAP3}

O sistema foi desenvolvido 	no Instituto Federal do Pará-Campus Castanhal, no laboratório do grupo de Computação Aplicada a Sistemas Inteligentes, durante os anos de 2015 e 2016.
\chapter{Resultados e Discussões}\label{CAP4}
\chapter{Conclusão}\label{CAP5}


%%%% Estilo de citação ABNT e arquivo de bibitens (mybibliography.bib)
\bibliographystyle{abnt-alf}
\bibliography{mybibliography}

\apendice
\chapter{Título do Apêndice}
\label{Apx:A}




\chapter{Exemplo do pacote Algorithm}
\label{Apx:B}


\begin{algorithm}[!h]
\caption{Estimador ML otimizado.}\label{Alg:MAXVER}
\begin{algorithmic}[1]
\STATE Inicializar o contador: $j\leftarrow 1$;%
\STATE Fixar o limiar de variação das estimativas: $e_{\mathrm{out}}\leftarrow 10^{-4}$;%
\STATE Fixar o número máximo de iterações: $N\leftarrow 1000$;%
\STATE Computar o ponto inicial: $\hat \gamma(0)$;%
\STATE Determinar o limiar inicial: $e_1 \leftarrow1000$;%
\STATE Estabelecer o valor inicial de $\alpha$: $\hat \alpha(0) \leftarrow -10^{-6}$;%
\WHILE{ $e_j \geq e_{\mathrm{out}}$ e $ j\leq M$}
    \STATE Solucionar $\hat \alpha_j\leftarrow {\arg \max}_{\alpha}\;{l_1(\alpha; \gamma_{j-1},\mathbf{z},n)}$;%
    \STATE Solucionar $\hat \gamma_j\leftarrow {\arg \max}_{\gamma}\;{l_2(\gamma; \alpha_j,\mathbf{z},n)}$;%
    \STATE $j\leftarrow j+1$
    \STATE Computar o critério de convergência: $e_j$;%
\ENDWHILE
\end{algorithmic}
\end{algorithm}


\end{document}